% ═══════════════════════════════════════════════════════════════════
% ПРИЛОЖЕНИЕ H: ПОЛНАЯ ТАБЛИЦА ВЫБОРКИ V2 И РЕЕСТР ФАЙЛОВ
% ═══════════════════════════════════════════════════════════════════

\section{Полная таблица выборки V2}\label{app:v2table}

Выборка V2 детерминирована: до трёх характеров на проводник
(первый, средний, последний в отсортированном списке), по три
данных обхода $(r,s)\in\{(0,0),(1,0),(0,1)\}$ на характер.
Ниже --- максимумы относительного отклонения по проводникам
(полные таблицы --- в JSON-отчётах \texttt{results/}).

\begin{center}
\small
\begin{tabular}{@{}l c c c@{}}
\toprule
Проводник & $N=15$: откл. & $N=30$: откл. & Порог PASS \\
\midrule
$d=3$  & $1.06\cdot10^{-36}$ & $1.06\cdot10^{-36}$ & $10^{-30}$ \\
$d=5$  & $2.88\cdot10^{-36}$ & $2.88\cdot10^{-36}$ & $10^{-30}$ \\
$d=6$  & ---                  & $2.68\cdot10^{-36}$ & $10^{-30}$ \\
$d=10$ & ---                  & $3.31\cdot10^{-36}$ & $10^{-30}$ \\
$d=15$ & $3.90\cdot10^{-36}$ & $3.90\cdot10^{-36}$ & $10^{-30}$ \\
$d=30$ & ---                  & $1.80\cdot10^{-36}$ & $10^{-30}$ \\
\midrule
Максимум & $3.90\cdot10^{-36}$ & $3.90\cdot10^{-36}$ & \\
Тестов & $21$ & $48$ & \\
\bottomrule
\end{tabular}
\end{center}

Отклонения на уровне $4\cdot10^{-36}$ согласованы с рабочей
точностью $dps=35$ (предел $\sim10^{-34}$): запас более чем на
порядок, и порог $10^{-30}$ проходит с резервом. Фазовые
отклонения V3 нулевые на всех $69$ тестах --- фаза вычисляется
в кольце вычетов без округлений.

\section{Реестр файлов лаборатории}\label{app:files}

\begin{center}
\small
\begin{tabularx}{\textwidth}{@{}l >{\raggedright\arraybackslash}X@{}}
\toprule
Файл / каталог & Содержание \\
\midrule
\texttt{laboratory.py} & единый файл-лаборатория: меню, протоколы,
отчёты, графики, мультиязычная верификация \\
\texttt{results/} & эталонные JSON-отчёты (базовая линия
сверки) \\
\texttt{reports/} & выходные отчёты прогонов (создаётся
лабораторией) \\
\texttt{verification/lean/} & \texttt{HodgeLaboratory.lean} ---
точные тождества в ядре Lean~4 \\
\texttt{verification/julia/} & \texttt{verify\_hodge.jl} ---
высокоточные периоды и фазы \\
\texttt{verification/fortran/} & \texttt{verify\_hodge.f90} ---
целочисленные переписи и ранги \\
\texttt{verification/c/} & \texttt{verify\_hodge.c} ---
быстрая целочисленная панель \\
\texttt{verification/rust/} & \texttt{verify\_hodge.rs} ---
безопасная арифметика \\
\texttt{monograph/} & PDF/DOCX монографии (RU, EN) + LaTeX-исходники \\
\texttt{theorems/} & шестнадцать папок --- отдельные монографии
теорем (RU, EN) \\
\texttt{docs/} & GitHub Pages: лендинг лаборатории \\
\texttt{scripts/} & установщик и автозагрузка репозитория
(Termux/Linux/macOS) \\
\texttt{LICENSE} & индивидуальная лицензия автора \\
\texttt{CITATION.cff} & метаданные цитирования \\
\bottomrule
\end{tabularx}
\end{center}

\section{Реестр проверок лаборатории}\label{app:labchecks}

Полный перечень проверок, встроенных в меню лаборатории, с
указанием раздела монографии и типа сертификата:

\begin{center}
\small
\begin{longtable}{@{}p{0.9cm} p{6.4cm} p{3.0cm} p{2.4cm}@{}}
\toprule
№ & Проверка & Раздел & Тип \\
\midrule
\endfirsthead
\toprule
№ & Проверка & Раздел & Тип \\
\midrule
\endhead
\bottomrule
\endfoot
T1 & перепись характеров (V1) & \ref{app:census} & точная \\
T2 & замкнутая форма vs tanh-sinh (V2) & \ref{sec:periods} &
интегрирование \\
T3 & фазы периода (V3) & \ref{sec:periods} & точная \\
T4 & эквивариантность (V4) & \ref{cor:equivar} & обе \\
T5 & целостность цепей (V5) & \ref{sec:protocol} &
интегрирование \\
T6 & лестница отражения (V6) & \ref{app:gamma} & обе \\
T7 & полнота/ранг (V7) & \ref{sec:protocol} & линейная алгебра \\
T8 & ортогональность ДПФ (V8) & \ref{sec:protocol} &
интегрирование \\
T9 & глубокая запись (V9) & \ref{sec:protocol} &
интегрирование \\
S1 & стенд тора (полный вывод) & \ref{sec:torus} & точная+float \\
S2 & стенд K3 ($\rho$, сигнатура, det) & \ref{sec:k3} & точная \\
S3 & errata E8 (перебор Арфа) & \ref{sec:e8} & точная \\
S4 & бинарный код (конвейер) & \ref{sec:binary} & точная \\
S5 & стенд Клейна ($j$, $\Omega$, $\tau$) & \ref{sec:klein} &
обе \\
C1--C8 & сертификаты A--H (ключевые проверки) &
\ref{sec:certA}--\ref{sec:certH} & обе \\
L1 & Lean-панель (ядро) & \ref{sec:multilang} & машинный вывод \\
X1--X5 & мультиязычная панель (Julia/Fortran/C/Rust) &
\ref{sec:multilang} & независимые реализации \\
E1 & конструктор собственных экспериментов &
\ref{sec:labmap} & параметризуемая \\
\end{longtable}
\end{center}

Суммарно лаборатория содержит $9$ проверок протокола, $5$ стендов,
$8$ сертификатов, $1$ Lean-панель, $5$ мультиязычных реализаций и
конструктор экспериментов --- все детерминированы, без сидов, с
автоматическими вердиктами PASS/FAIL и генерацией отчётов.
